\documentclass[11pt]{article}

\DeclareUnicodeCharacter{0301}{\'e}

\usepackage[margin=1in]{geometry}
\usepackage{multirow}
\usepackage{amssymb,amsmath,amsthm,amsfonts}
\usepackage{bm}
\usepackage{textgreek}
\usepackage{mathtools}
\usepackage{enumitem}
\usepackage[numbers,comma,sort&compress]{natbib}
\usepackage{authblk}
\usepackage{graphicx}
\usepackage[font=small]{caption}
\usepackage{subcaption} % [labelformat=simple]
\usepackage{tablefootnote} 
\usepackage{float}
\usepackage[ruled,vlined,linesnumbered]{algorithm2e}
\usepackage{algorithmic}
\renewcommand{\algorithmicrequire}{\textbf{Input: }}  
\renewcommand{\algorithmicensure}{\textbf{Output: }}
\usepackage{physics}
\usepackage{footnote}
\usepackage{xcolor}
\usepackage{mathrsfs}
\usepackage{bbm}
\usepackage{makecell}
\usepackage{pbox}
\usepackage[colorlinks]{hyperref}
\hypersetup{citecolor=blue}
% \urlstyle{same}
%\captionsetup[figure]{labelfont=bf}

\newtheorem{fact}{Fact}[section] 
\newtheorem{theorem}{Theorem}
\newtheorem{definition}{Definition}
\newtheorem{lemma}{Lemma}
\newtheorem{proposition}{Proposition}
\newtheorem{example}{Example}
\newtheorem{corollary}{Corollary}
\newtheorem{remark}{Remark}
\newtheorem{assumption}{Assumption}
\newcommand{\eqn}[1]{(\ref{eqn:#1})}
\newcommand{\eq}[1]{(\ref{eq:#1})}
\newcommand{\prb}[1]{(\ref{prb:#1})}
\newcommand{\thm}[1]{\hyperref[thm:#1]{Theorem~\ref*{thm:#1}}}
\newcommand{\cor}[1]{\hyperref[cor:#1]{Corollary~\ref*{cor:#1}}}
\newcommand{\defn}[1]{\hyperref[defn:#1]{Definition~\ref*{defn:#1}}}
\newcommand{\lem}[1]{\hyperref[lem:#1]{Lemma~\ref*{lem:#1}}}
\newcommand{\prop}[1]{\hyperref[prop:#1]{Proposition~\ref*{prop:#1}}}
\newcommand{\assum}[1]{\hyperref[assum:#1]{Assumption~\ref*{assum:#1}}}
\newcommand{\fig}[1]{\hyperref[fig:#1]{Figure~\ref*{fig:#1}}}
\newcommand{\tab}[1]{\hyperref[tab:#1]{Table~\ref*{tab:#1}}}
\newcommand{\algo}[1]{\hyperref[algo:#1]{Algorithm~\ref*{algo:#1}}}
\renewcommand{\sec}[1]{\hyperref[sec:#1]{Section~\ref*{sec:#1}}}
\newcommand{\append}[1]{\hyperref[append:#1]{Appendix~\ref*{append:#1}}}
\newcommand{\fac}[1]{\hyperref[fac:#1]{Fact~\ref*{fac:#1}}}
\newcommand{\lin}[1]{\hyperref[lin:#1]{Line~\ref*{lin:#1}}}
\newcommand{\itm}[1]{\ref{itm:#1}}
\newcommand{\bdelta}{\mathbf{\Delta}}
\renewcommand{\arraystretch}{1.25}
\newcommand{\algname}[1]{\textup{\texttt{#1}}}
\renewcommand{\i}{\mathrm{i}}
\renewcommand{\d}{\mathrm{d}}
\renewcommand{\L}{\mathcal{L}}
\newcommand{\ad}{\mathrm{ad}}
\newcommand{\T}{\mathcal{T}}
\newcommand{\B}{\Big}
\renewcommand{\P}{\mathcal{P}}
\newcommand{\simu}{\mathrm{sim}}
\newcommand{\Hc}{H_{\mathrm{control}}}
\newcommand{\He}{H_{\mathrm{eff}}}
\def\>{\rangle}
\def\<{\langle}
\def\trans{^{\top}}

\newcommand{\xz}[1]{{\color{cyan} XZ: #1}}
\newcommand{\tyl}[1]{{\color{blue} Tongyang: #1}}
\newcommand{\zsy}[1]{{\color{teal} Shengyu: #1}}
\newcommand{\sz}[1]{{\color{purple} Shuo: #1}}



\title{Lindbladian Simulation with ZNE}
\author{}
\date{}

\begin{document}

\maketitle
\section{Error Bounds of the Richardson Extrapolation}
Given an $n$-qubit Hamiltonian $H = \sum_{\gamma=1}^{\Gamma} H_{\gamma}$, we aim to simulate $e^{-\i H T}$ under Markovian physical noise. Let $\P_2(t)=e^{-\i H_1 t/2}\cdots e^{-\i H_{\Gamma} t/2} e^{-\i H_{\Gamma} t/2} \cdots e^{-\i H_1 t/2}$ be the second order Trotter formula. The simulation is performed by repeatedly applying $\P_2(sT)$ for $1/s$ times. \textbf{Here we assume that to implement the Trotter step for different step size, we only need to linearly scale the control Hamiltonian of the physical device.} Specifically, each application of $\P_2(sT)$ is realized by evolving the quantum computer under a time-dependent control Hamiltonian $\Hc(t/s)$ for $t\in [0, sT_{\simu}]$, where $T_{\simu}$ is the physical time required to execute one Trotter step of simulated time $T$. Formally, we have \begin{align*}
    \P_2(sT) = \T \exp(\int_{0}^{sT_{\simu}} - \i \Hc(t/s) \, \d t).
\end{align*}
Concurrently, the physical device is subject to Markovian noise described by a Lindbladian $\lambda \L$. While implementing one Trotter step $\P_2(sT)$ of step size $sT$, the dynamics of the physical device's density matrix $\tilde{\rho}_s(t)$ is described by \begin{align*}
    \frac{\d \tilde{\rho}_s(t)}{\d t} = -\i\ad_{\Hc(t/s)}(\tilde{\rho}_s(t))+\lambda \L(\tilde{\rho}_s(t)),
\end{align*}
for $t\in[0, sT_{\simu}]$, where $\ad_H(\rho) :=[H,\rho]$.
We define the ideal time-evolution operator and the time-evolution operator with noise strength $\lambda$ as \begin{align*}
    \Phi_s(t_2,t_1):=\T \exp(\int_{t_1}^{t_2} -\i\ad_{\Hc(t/s)}\,\d t),\quad \tilde{\Phi}_{s,\lambda}(t_2, t_1):=\T \exp(\int_{t_1}^{t_2}(-\i \ad_{\Hc(t/s)}+\lambda \L)\,\d t). 
\end{align*}
\subsection{Method 1: the Variation of Parameter Formula}
By applying the formula variation of parameters $K$ times, the noisy time evolution operator, $\tilde{\Phi}_{s,\lambda}$, can be expanded as a series of powers in the noise strength $\lambda$:
\begin{align}
\label{eq:expansion-s-l}
    \tilde{\Phi}_{s,\lambda}(s T_{\simu}, 0) = \Phi_s(s T_{\simu}, 0) + \sum_{\ell=1}^{K-1}F_{s,\ell}\lambda^{\ell} + F_{s,K}(\lambda)\lambda^K 
\end{align}
where the expansion terms are given by the time-ordered integrals:
\begin{align}
    \label{eq:def-Fsl}
    F_{s,\ell} := \int_{0}^{s T_{\simu}} \d \tau_1 \int_{0}^{\tau_1}\d \tau_2 \cdots \int_{0}^{\tau_{\ell-1}} \d \tau_{\ell} \Phi_s(s T_{\simu}, \tau_1) \L \Phi_s(\tau_1, \tau_2) \L \cdots \Phi_s(\tau_{\ell-1}, \tau_{\ell}) \L\Phi_s(\tau_{\ell},0),
\end{align}
and
\begin{align*}
    F_{s, K}(\lambda) := \int_{0}^{s T_{\simu}} \d \tau_1 \int_{0}^{\tau_1}\d \tau_2 \cdots \int_{0}^{\tau_{K-1}} \d \tau_{K} \Phi_s(s T_{\simu}, \tau_1) \L \Phi_s(\tau_1, \tau_2) \L \cdots \Phi_s(\tau_{K-1}, \tau_{K}) \L\tilde{\Phi}_{s,\lambda}(\tau_{K},0).
\end{align*}
Substituting $\tau_j=s\sigma_j$ into \eq{def-Fsl} and setting $\sigma_{0}:=T_{\simu}, \sigma_{\ell+1}:=0$, we get \begin{align}
\label{eq:Fsl}
    F_{s,\ell}=s^\ell \int_{0}^{T_{\simu}} \d\sigma_1 \cdots \int_{0}^{\sigma_{\ell-1}} \d\sigma_{\ell} \B( \prod_{j=1}^{\ell} \Phi_s(s\sigma_{j-1}, s\sigma_j) \L \B) \Phi_s(s\sigma_\ell, s\sigma_{\ell+1}) ,
\end{align}
where \begin{align*}
    \Phi_s(s\sigma_{j-1}, s\sigma_j) &=\T \exp(\int_{s\sigma_j}^{s\sigma_{j-1}} -\i\ad_{\Hc(t/s)}\,\d t) \\
    &= \T \exp(s\int_{\sigma_j}^{\sigma_{j-1}} -\i\ad_{\Hc(t)}\,\d t) \\
\end{align*}
has Dyson series expansion \begin{align}
\label{eq:dyson-Phi}
    \Phi_s(s\sigma_{j-1}, s\sigma_j)=\sum_{k=0}^{\infty} s^k V_k,\quad  \|V_k\| \leq \frac{(2 H_{\max }(\sigma_{j-1}-\sigma_j))^k}{k!}.
\end{align}
Substituting Eq.~\eq{dyson-Phi} into Eq.~\eq{Fsl}, we get \begin{align*}
    F_{s,\ell} = \sum_{p = \ell}^{\infty} s^p C_{p,\ell},
\end{align*}
where \begin{align}
    \|C_{p,\ell}\| &\le \|\L\|^{\ell} \int_{0}^{T_{\simu}}\d\sigma_1 \cdots \int_{0}^{\sigma_{\ell-1}} \d\sigma_{\ell} \sum_{k_1+\dots+k_{\ell+1}=p-\ell} \prod_{j=1}^{\ell+1} \frac{(2 H_{\max }(\sigma_{j-1}-\sigma_j))^{k_j}}{k_j!}\\
     &= \|\L\|^{\ell}\int_{0}^{T_{\simu}}\d\sigma_1 \cdots \int_{0}^{\sigma_{\ell-1}} \d\sigma_{\ell} \frac{(\sum_{j = 1}^{\ell+1}2 H_{\max }(\sigma_{j-1}-\sigma_j))^{p-\ell}}{(p-\ell)!}\\
    &= \|\L\|^{\ell} \int_{0}^{T_{\simu}}\d\sigma_1 \cdots \int_{0}^{\sigma_{\ell-1}} \d\sigma_{\ell} \frac{(2H_{\max}T_{\simu})^{p-\ell}}{(p-\ell)!}\\
    & =  \frac{T_{\simu}^p \|\L\|^\ell (2H_{\max})^{p-\ell}}{\ell!(p-\ell)!}.
\end{align}
Taking $K\to \infty$ and substituting Eq.~\eq{dyson-Phi} into Eq.~\eq{expansion-s-l}, we get 
\begin{align}
\label{eq:expansion-s}
    \tilde{\Phi}_{s,\lambda}(s T_{\simu}, 0) = \Phi_s(s T_{\simu}, 0)+\sum_{\ell=1}^{\infty}\sum_{p\ge \ell}^{\infty}s^p \lambda^{\ell} C_{p,\ell} 
\end{align}
\subsection{Method 2: the Magnus Expansion}
\begin{theorem}[{\cite[Theorem III]{magnus1954}}]
\label{thm:magnus-exp}
    Let $A(t)$ be a function of $t$ in an associative ring, and let \begin{align*}
        Y(t) = \T \exp(\int_{0}^{t} A(s)\,\d s).
    \end{align*}
    Then, if certain unspecified conditions of convergence are satisfied, $Y(t)$ can be written in the form \begin{align*}
        Y(t) = \exp \Omega(t),
    \end{align*}
    where 
    \begin{align*}
        \frac{\d \Omega}{\d t}=\sum_{n=0}^{\infty} \frac{B_n}{n!} \ad_{\Omega}^n A
    \end{align*}
    and $B_n$ are the Bernoulli numbers.
\end{theorem}
Take $A_{s,\lambda}(t) := s(-\i \ad_{\Hc(t)}+\lambda\L)$ in \thm{magnus-exp}, and then we have \begin{align}
    Y_{s,\lambda}(T_{\simu}) & = \T\exp(\int_{0}^{T_\simu} s(-\i \ad_{\Hc(t)}+\lambda\L)\,\d t)\\
    & = \T\exp(\int_{0}^{sT_\simu}(-\i \ad_{\Hc(t)}+\lambda\L)\,\d \tau) \\
    & =\tilde{\Phi}_{s,\lambda}(sT_{\simu}, 0).
\end{align}
By \thm{magnus-exp}, if the convergence conditions are satisfied, $\tilde{\Phi}_{s,\lambda}(sT_{\simu}, 0)$ can be written as \begin{align*}
    \tilde{\Phi}_{s,\lambda}(sT_{\simu}, 0) = \exp \Omega_{s,\lambda}(T_{\simu}).
\end{align*}
Following the proof of \cite[Theorem 6]{blanes2009magnus}, we can show that $\Omega_{s,\lambda}(T_{\simu})$ is an analytic function of $s$ and $\lambda$.
\begin{lemma}
    \label{lem:analytic-phi}
    The function $\varphi(s,\lambda):=\log(Y_{s,\lambda}(T_{\simu}))$ is an analytic function of $s$ and $\lambda$ in the set $B_{\pi}$, with \begin{align*}
        B_{\pi} :=\{s,\lambda \in \mathbb{C}: \int_{0}^{T_{\simu}}\|A_{s,\lambda}(t) \|\,\d t\le \pi \}.
    \end{align*}
\end{lemma}
 
Now we take $\lambda = Cs$. 
\begin{lemma}
\label{lem:remainder-f}
    Given $C> 0$, let $\He(s) := \frac{\varphi(s, Cs)}{sT_{\simu}}$ and \begin{align*}
        f(s):= \exp(\He(s) T_{\simu}).
    \end{align*}
    Let \begin{align*}
        r=\min\Big\{ \frac{1}{4\int_{0}^{T_{\simu}}\|\Hc(t)\|\, \d t}, \frac{1}{2\sqrt{C\|\L\|T_{\simu}}}\Big\}.
    \end{align*}
    For any integer $k\ge 1$, there exists a degree-$k$ polynomial $P_k(s)\colon \mathbb{C}\to \mathbb{C}^{2^n\times 2^n}$ and a remainder $R_k(s)\colon \mathbb{C}\to \mathbb{C}^{2^n\times 2^n}$, such that \begin{align*}
        f(s) = P_k(s)+R_k(s), \text{ and }
        \|R_k(s)\|=O\B(\frac{|s|^{k+1}}{r^k(r-|s|)}\B)
    \end{align*}
    for all $|s|\le \min\{r/2,r^2/(8\pi)\}$. 
\end{lemma}
\begin{proof}
    For $|s|< r$, we have \begin{align}
    \label{eq:converge-radius}
        \int_{0}^{T_{\simu}}\|A_{s,Cs}(t) \|\,\d t\le\int_{0}^{T_{\simu}}(2|s|\|\Hc(t)\|+C|s|\|\L\|)\, \d t\le 1.
    \end{align}
    Let $g(s) = \varphi(s, Cs)$. By \lem{analytic-phi}, $g(s)$ is a holomorphic function of $s$ for $|s|< r$. Since $\lim_{s\to 0}g(s)= 0$, 
    $\He(s)=\frac{g(s)}{sT_{\simu}}$ and $f(s) = \exp(\He(s)T_{\simu})$ are holomorphic functions of $s$ for $|s|< r$. For $s=0$, we define $\He(0) = \lim_{s\to 0} \He(s) = g'(0)/T_{\simu}= -\i HT/T_{\simu}$. \xz{TBD by comparing the series expansion of $\exp(\varphi(s, Cs))$ and noisy $\P_2(sT)$.}
    \citet{moan1998efficient} showed that $\|\varphi(s,C s)\|\le \pi$ for  $ \int_{0}^{T_{\simu}}\|A_{s, Cs}(t) \|\,\d t\le 1.08686871$. From Eq.~\eq{converge-radius}, the norm $\|g(s)\|\le \pi $ for all $|s|\le r$. 

    
    In the interaction picture, $f(s)$ can be written as \begin{align*}
        f(s)&=\exp(\He(s) T_{\simu}) \\
        &= \exp(\He(0)T_{\simu})\B(I+\T\exp\B(\int_{0}^{T_{\simu}}e^{-t\He(0)} (\He(s)-\He(0))e^{t\He(0)}\, \d t\B)\B).
    \end{align*}
    Note that $\He(0)=-\i HT/T_{\simu}$ is anti-hermitian, we can bound the norm of $f(s)$ as \begin{align*}
        \|f(s)\|&\le \|\exp(\He(0)T_{\simu})\|\B(1+\B\|\T\exp\B(\int_{0}^{T_{\simu}}e^{-t\He(0)} (\He(s)-\He(0))e^{t\He(0)}\, \d t\B)\B\|\B)\\
        &\le 1+\exp\B(\int_{0}^{T_{\simu}}\B\|e^{-t\He(0)} (\He(s)-\He(0))e^{t\He(0)}\B\|\, \d t\B)\\
        &=1+\exp(\|\He(s)-\He(0)\|T_{\simu}).
    \end{align*}
    To bound the exponent, we first give a bound for $\|g''(s)\|$: 
    \begin{align*}
        \|g''(s)\| = \B\|\frac{1}{2\pi i}\int_{|z|=r} g(z) \frac{2}{(z-s)^3}\, \d z\B\| \le \max_{|z|\le r}\|g(z)\| \frac{2r}{(r-|s|)^3}\le \frac{16\pi}{r^2} ,
    \end{align*}
    for $|s|\le r/2$.
    The exponent $\|\He(s)-\He(0)\|T_{\simu}$ can be bounded as \begin{align*}
        \|(\He(s)-\He(0)\|T_{\simu} &= \frac{1}{|s|}\|g(s)-g'(0)s\|\\
        &=\frac{1}{|s|}\B\|\int_0^s g''(z)(s-z)\, \d z\B\| \\
        &\le \frac{1}{|s|}\max_{|z|\le s}\|g''(z)\| \frac{|s|^2}{2} \\
        &\le \frac{|s|}{2}\frac{16\pi}{r^2} \le 1,
    \end{align*}
    for $|s|\le \min\{r/2,r^2/(8\pi)\}$. Then $\|f(s)\|\le 1+e$ for $s$ in the same region. Then for any integer $k\ge 1$, $f(s)$ can be decomposed as \begin{align*}
        f(s) = P_k(s)+R_k(s),
    \end{align*}
    where $P_k(s)$ is the $k$-th-order Taylor polynomial and $R_k(s)$ is the remainder satisfying \begin{align*}
        \|R(s)\| &=\B\|\frac{s^{k+1}}{2\pi \i}\int_{|z|=r}\frac{f(z)}{z^{k+1}(z-s)}\, \d z\B\|\\
        &\le \max_{|z|\le r}\|f(z)\| \frac{|s|^{k+1}}{r^{k}(r-|s|)}\le (1+e)\frac{|s|^{k+1}}{r^{k}(r-|s|)}.
    \end{align*}
    \xz{Here I directly use some results of $\mathbb{C}\to \mathbb{C}$ holomorphic functions for $\mathbb{C}\to \mathbb{C}^{2^n \times 2^n}$ holomorphic functions, which need to further checked the correctness.}
\end{proof}
The scaling of the Trotter step size in \lem{remainder-f} is similar to qDRIFT, since the first-order noise dominants. Intuitively, if we take the noise strength $\lambda = Cs^{p}$ and use $p$-th-order product formulas, $s=O(r^{1+1/p})$ suffices. The key is to prove that the $g^{(j)}(0)=0$ for $2\le j\le p+1$, where $g(s) = \varphi(s, Cs^p)$.

Take $\lambda = C s^p$ and $\P$ be a $p$-th-order product formula. By the variation of parameter formula, we have \begin{align*}
    Y_{s,Cs^p}(T_{\simu}) & = \T \exp(s\int_0^{T_{\simu}}-\i \ad_{\Hc(t)} \, \d t)\\
    &\quad+\int_{0}^{T_{\simu}} \exp_{\T}(s\int_{t}^{T_{\simu}}-\i \ad_{\Hc(\tau)}\, \d \tau )Cs^{p+1}\L\exp_{\T}(s\int_{0}^{t}(-\i \ad_{\Hc(\tau)}+Cs^p\L)\, \d \tau )\, \d t \\
     &= \P(sT)+O(s^{p+1}) = e^{-\i HsT}+O(s^{p+1})
\end{align*}
Then we have \begin{align*}
    \varphi(s,Cs^{p}) &= \log(Y_{s,Cs^p}(T_{\simu}))\\
    & =\sum_{k=1}^{\infty}\frac{(-1)^{k-1}}{k}(Y_{s,Cs^p}(T_{\simu})-I)^k \\
    & = \log(e^{-\i H s T})+O(s^{p+1})=-\i HsT + O(s^{p+1}).
\end{align*}

\bibliographystyle{plainnat}
\bibliography{ref}
\end{document}
